% Hauptdatei der Bachelorarbeit
% Kompilieren mit: latexmk -pdf thesis.tex
\documentclass[
  a4paper,
  11pt,
  oneside,            % Keine Ausgleichsseiten für einen zweiseitigen Druck.
  openany,
  parskip=half,
  headings=normal,
  numbers=noenddot,
  bibliography=totoc,
  listof=totoc,
  captions=tableheading,
  cleardoublepage=empty
]{scrbook}

% Aktueller Titel: Automatisierte Informationsextraktion aus ausgefüllten Teilnehmerlisten: Analyse KI-basierter Ansätze hinsichtlich Qualität, Effizienz und praktischer Einsatzfähigkeit
% Englischer Titel: Automated Information Extraction from Participant Lists: Analysis of AI-Based Approaches in Terms of Quality, Efficiency, and Practical Applicability
% Titelbestandteile werden für das Titelblatt in config/metadata.tex gepflegt.
\newcommand{\researchquestion}{Wie unterscheiden sich verschiedene \acs{ki}-basierte Ansätze zur automatisierten Informationsextraktion aus handschriftlich ausgefüllten Teilnehmerlisten hinsichtlich Extraktionsqualität, Effizienz und praktischer Einsatzfähigkeit?}

% !TeX root = ../thesis.tex
%%%%%%%%%%%%%%%%%%%%%%%%%%%%%%%%%%%%%%%%%%%%%%%%%%%%%%%%%%%%%%%%%%%%%%%%%%%%%%%%
% Zentrale Angaben zur Arbeit
% Nur diese Datei muss für Titelblatt und PDF-Metadaten angepasst werden.
%%%%%%%%%%%%%%%%%%%%%%%%%%%%%%%%%%%%%%%%%%%%%%%%%%%%%%%%%%%%%%%%%%%%%%%%%%%%%%%%

\newcommand{\thesismaintitleDE}{Automatisierte Informationsextraktion aus handschriftlich ausgefüllten Teilnehmerlisten}
\newcommand{\thesissubtitleDE}{Analyse \texorpdfstring{\acs{ki}}{KI}-basierter Ansätze hinsichtlich Qualität, Effizienz und praktischer Einsatzfähigkeit}
% Vollständiger Titel für PDF-Metadaten und andere Verwendungen.
\newcommand{\thesistitleDE}{\thesismaintitleDE: \thesissubtitleDE}
\newcommand{\thesistitleEN}{Automated Information Extraction from Handwritten Participant Lists: Analysis of AI-Based Approaches in Terms of Quality, Efficiency, and Practical Applicability}
\newcommand{\thesistype}{Bachelorarbeit}

\newcommand{\student}{Marco Martinetz}
\newcommand{\matrnr}{22204884}
\newcommand{\stammnr}{} % Leer lassen, falls nicht benötigt.
\newcommand{\submissiondate}{15.03.2027}

\newcommand{\supervisor}{Prof. Dr. Robert Hable}
\newcommand{\secsupervisor}{} % Leer lassen, falls nicht benötigt.

\newcommand{\faculty}{Angewandte Informatik}
\newcommand{\studies}{Bachelor Künstliche Intelligenz}
\newcommand{\degree}{Bachelor of Science (B.Sc.)}

% TODO-Hinweise im PDF anzeigen. Vor der finalen Abgabe auf \showtodosfalse setzen
% und zusätzlich das gesamte Projekt nach "TODO" durchsuchen.
\newif\ifshowtodos
\showtodostrue

% !TeX root = ../thesis.tex
% Kodierung, Sprache und Schrift
\usepackage[utf8]{inputenc}
\usepackage[T1]{fontenc}
\usepackage[english,ngerman]{babel}
\usepackage{lmodern}
\usepackage{microtype}
\usepackage{csquotes}
\usepackage{acronym}
% Der Titel wird vor dem Abkürzungsverzeichnis gesetzt.
\newacro{ki}[KI]{Künstliche Intelligenz}

% Mathematik, Tabellen und Abbildungen
\usepackage{amsmath,amssymb}
\usepackage{booktabs}
\usepackage{tabularx}
\usepackage{graphicx}
\usepackage{xcolor}
\graphicspath{{figures/}}

% Hilfspakete
\usepackage{ifthen}
\usepackage[automark,headsepline]{scrlayer-scrpage}
\clearpairofpagestyles
\ihead{\headmark}
% Sternform setzt die Seitenzahl auch auf Kapitel- und Verzeichnisanfangsseiten.
\cfoot*{\pagemark}
\pagestyle{scrheadings}

% Hyperlinks möglichst zuletzt laden.
\usepackage{hyperref}
\hypersetup{
  hidelinks,
  pdftitle={\thesistitleDE},
  pdfauthor={\student},
  pdfsubject={\thesistype},
  pdfcreator={LaTeX}
}
\urlstyle{same}

% Satzspiegel und Gliederungstiefe
\KOMAoptions{DIV=12,BCOR=8mm}
\setcounter{secnumdepth}{3}
\setcounter{tocdepth}{3}

% Geplante Inhalte der vorläufigen Gliederung.
\newenvironment{inhalte}{%
  \begin{list}{\textbullet}{%
    \setlength{\leftmargin}{1.45em}%
    \setlength{\labelwidth}{.7em}%
    \setlength{\labelsep}{.5em}%
    \setlength{\topsep}{.15em}%
    \setlength{\itemsep}{0pt}%
    \setlength{\parsep}{0pt}%
    \setlength{\partopsep}{0pt}%
  }%
}{\end{list}}

% Sichtbare Platzhalter für die Schreibphase.
\newcommand{\TODO}[1]{%
  \ifshowtodos
    \par\smallskip
    \noindent{\color{red!70!black}\textbf{TODO:} #1}
    \par\smallskip
  \fi
}


\begin{document}

\frontmatter
% !TeX root = ../thesis.tex
\begin{titlepage}
  \centering

  \IfFileExists{assets/THD_Logo.pdf}{%
    \includegraphics[width=5.2cm]{assets/THD_Logo.pdf}\par
    \vspace{1.2cm}
  }{}

  {\Large Technische Hochschule Deggendorf\par}
  \vspace{0.25cm}
  {\large Fakultät \faculty\par}
  {\large Studiengang \studies\par}

  \vfill

  {\Huge\bfseries \thesismaintitleDE\par}
  \vspace{0.7cm}
  {\Large \thesissubtitleDE\par}
  \ifthenelse{\equal{\thesistitleEN}{}}{}{%
    \vspace{0.8cm}
    {\Large \thesistitleEN\par}
  }

  \vfill

  {\Large \thesistype\par}
  \vspace{0.4cm}
  zur Erlangung des akademischen Grades
  \vspace{0.25cm}

  {\large\itshape \degree\par}

  \vfill

  \begin{minipage}[t]{0.46\textwidth}
    \raggedright
    Vorgelegt von: \student\\

    Matrikelnummer: \matrnr
    \ifthenelse{\equal{\stammnr}{}}{}{\\Stammnummer: \stammnr}

    \vspace{\baselineskip}
    Abgabedatum: \submissiondate
  \end{minipage}
  \hfill
  \begin{minipage}[t]{0.46\textwidth}
    \raggedright
    Erstprüfer/in: \supervisor

    \ifthenelse{\equal{\secsupervisor}{}}{}{%
      \vspace{\baselineskip}

      Zweitprüfer/in:\\
      \secsupervisor
    }
  \end{minipage}
\end{titlepage}
\cleardoublepage

% !TeX root = ../thesis.tex
\chapter{Kurzfassung}

\TODO{Die Kurzfassung auf Deutsch verfassen. Sie sollte Problemstellung, Ziel, Vorgehen, wichtigste Ergebnisse und Schlussfolgerung knapp zusammenfassen.}

\textbf{Schlagwörter:} \TODO{Drei bis fünf Schlagwörter ergänzen.}

% !TeX root = ../thesis.tex
\begin{otherlanguage}{english}
\chapter{Abstract}

\TODO{Write the English abstract. It should briefly cover the problem, objective, method, principal results, and conclusion.}

\textbf{Keywords:} \TODO{Add three to five keywords.}
\end{otherlanguage}


\tableofcontents
% Aktivieren, sobald Abbildungen bzw. Tabellen vorhanden sind.
% \listoffigures
% \listoftables
% !TeX root = ../thesis.tex
\chapter{Abkürzungsverzeichnis}

% Alphabetisch sortiert; neue Abkürzungen hier definieren und im Text mit \ac markieren.
\begin{acronym}[GNHK]
  \acro{api}[API]{Programmierschnittstelle (Application Programming Interface)}
  \acro{cer}[CER]{Character Error Rate}
  \acro{ctc}[CTC]{Connectionist Temporal Classification}
  \acro{dia}[DIA]{Document Image Analysis}
  \acro{dla}[DLA]{Document Layout Analysis}
  \acro{erp}[ERP]{Enterprise Resource Planning}
  \acro{gnhk}[GNHK]{Bezeichnung des Handschriftdatensatzes von Lee et al.}
  \acro{htr}[HTR]{Handwritten Text Recognition}
  \acro{json}[JSON]{JavaScript Object Notation}
  \acro{ki}[KI]{Künstliche Intelligenz}
  \acro{llava}[LLaVA]{Large Language and Vision Assistant}
  \acro{llm}[LLM]{Large Language Model}
  \acro{ocr}[OCR]{Optical Character Recognition}
  \acro{pdf}[PDF]{Portable Document Format}
  \acro{trocr}[TrOCR]{Transformer-based Optical Character Recognition}
  \acro{vlm}[VLM]{Vision-Language Model}
\end{acronym}


\mainmatter
% !TeX root = ../thesis.tex
% Kapitel 1: Entwurf, zuletzt am 08.10.2026 überarbeitet; noch nicht final abgestimmt.
\chapter{Einleitung}
\label{chap:einleitung}

\section{Problemstellung und Motivation}

Im betrachteten Unternehmensprozess wird die Teilnahme an Schulungen anhand handschriftlich ausgefüllter Teilnehmerlisten dokumentiert. Die darin enthaltenen Angaben werden bislang durch eine Person ausgelesen und in ein betriebliches Schulungssystem übertragen. Zu den für die Untersuchung vorgesehenen Informationen gehören Vorname, Nachname, Personalnummer, Abteilung und Angaben zum Vorhandensein einer Unterschrift. Die manuelle Übertragung bildet den Ausgangspunkt für die Frage, in welchem Umfang \ac{ki} die Erfassung unterstützen kann und welcher Aufwand für eine verlässliche Weiterverarbeitung verbleibt.

Die für die Untersuchung vorgesehenen historischen Dokumente liegen als Dateien im \ac{pdf} vor. Diese können neben handschriftlich ausgefüllten Teilnehmerlisten auch Titelseiten, leere Seiten, nicht ausgefüllte oder maschinell ausgefüllte Listen sowie weitere Inhalte enthalten. Hinzu kommen unterschiedliche Listenvarianten und abweichende Seitenorientierungen. Vor der eigentlichen Extraktion müssen deshalb die relevanten Seiten identifiziert und in eine geeignete Orientierung gebracht werden. Diese Eigenschaften der Datengrundlage verdeutlichen, dass die Verarbeitung bereits vor der Erkennung einzelner Feldwerte mehrere aufeinander aufbauende Aufgaben umfasst.

Für die Erkennung handschriftlicher Inhalte ist zudem die Beschaffenheit des Bildmaterials relevant. Lee et al.~\cite{Lee2021GNHK} untersuchen mit dem \acs{gnhk}-Datensatz die Lokalisierung und Erkennung englischer Handschriften in fotografierten Dokumenten aus unterschiedlichen realen Umgebungen. Die Arbeit verdeutlicht die Bedeutung von Untersuchungsbedingungen, die über standardisierte Textausschnitte hinausgehen. Für die vorliegende Untersuchung folgt daraus, dass die Dokumenteigenschaften des betrieblichen Anwendungsfalls bei der Zusammenstellung und Auswertung der Testdaten berücksichtigt werden müssen.

Die angestrebte Automatisierung geht über die Überführung eines Schriftbilds in eine Zeichenfolge hinaus. Erkannte Angaben müssen den vorgesehenen Feldern und den richtigen Teilnehmenden zugeordnet werden. Ein korrekt gelesener Name ist für die weitere Verarbeitung nur dann nutzbar, wenn auch seine Zuordnung zur zugehörigen Personalnummer und zu den weiteren Angaben stimmt. Skalický et al.~\cite{Skalicky2022Benchmarks} unterscheiden entsprechend zwischen der Extraktion einzelner Schlüsselinformationen und der Erkennung zusammengehöriger Einträge innerhalb eines Dokuments. Sie heben außerdem die Bedeutung der Lokalisierung extrahierter Angaben für Anwendungen mit menschlicher Validierung hervor. Diese Unterscheidungen sind auf die Bewertung mehrzeiliger Teilnehmerlisten übertragbar.

Zur Bearbeitung solcher Aufgaben stehen unterschiedliche \acs{ki}-basierte Ansätze zur Verfügung. \ac{trocr} zeigt einen auf vortrainierten Transformern beruhenden Ansatz zur Texterkennung, der unter anderem auf handschriftlichen Textzeilen evaluiert wird~\cite{Li2023TrOCR}. Donut verfolgt dagegen die direkte Verarbeitung von Dokumentbildern für Aufgaben des Dokumentenverständnisses, ohne dass vorher \ac{ocr} angewandt wird. Die Autoren diskutieren dabei unter anderem die Weitergabe von \acs{ocr}-Fehlern an nachgelagerte Verarbeitungsschritte sowie den Aufwand mehrstufiger Pipelines~\cite{Kim2022Donut}. Diese Arbeiten illustrieren unterschiedliche Verarbeitungsarchitekturen, deren Eignung für den vorliegenden Anwendungsfall empirisch zu prüfen ist.

Auch die Fähigkeiten großer multimodaler Modelle begründen keine voraussetzungslose Zuverlässigkeit bei der Dokumentenextraktion. OCRBench v2 dokumentiert bei den untersuchten Modellen Grenzen in textbezogenen Bildaufgaben, darunter Handschriftextraktion, Textlokalisierung und Layoutverständnis~\cite{Fu2025OCRBenchV2}. Daraus lässt sich keine konkrete Fehlerquote für die hier betrachteten Teilnehmerlisten ableiten. Die Ergebnisse begründen jedoch die Notwendigkeit, ausgewählte Verfahren anhand einer zum Anwendungsfall passenden Datengrundlage zu überprüfen, anstatt allgemeine Modellfähigkeiten mit betrieblicher Eignung gleichzusetzen.

Neben der Extraktionsqualität ist zu untersuchen, welcher Nutzen im bestehenden Arbeitsablauf erreichbar ist. Eine vergleichende Studie zur Digitalisierung gedruckter Berichte aus dem Bereich \ac{erp} verbindet beispielsweise die Bewertung multimodaler Modelle mit einer Untersuchung manueller Datenerfassung hinsichtlich Genauigkeit, Zeit und Kosten~\cite{Lertyosbordin2026ERP}. Ihre Ergebnisse sind aufgrund der abweichenden Dokumente und Prozessbedingungen nicht unmittelbar auf handschriftliche Teilnehmerlisten übertragbar. Für die vorliegende Arbeit werden deshalb eigene Prozessdaten benötigt. Mögliche Einsparungen müssen unter Berücksichtigung der automatisierten Verarbeitung, der menschlichen Kontrolle und Korrektur sowie der weiterhin notwendigen Übernahme in das Schulungssystem beurteilt werden.

\section{Zielsetzung der Arbeit}

Ziel dieser Arbeit ist die vergleichende Untersuchung ausgewählter \acs{ki}-basierter Ansätze zur automatisierten Informationsextraktion aus handschriftlich ausgefüllten Teilnehmerlisten. Im Mittelpunkt stehen die korrekte und vollständige Erfassung von Vorname, Nachname, Personalnummer und Abteilung sowie die Feststellung, ob im vorgesehenen Feld eine Unterschrift vorhanden ist. Die Angaben sollen strukturiert ausgegeben und den jeweiligen Teilnehmenden innerhalb der Liste zugeordnet werden. Die Bewertung des Unterschriftsfelds beschränkt sich ausschließlich auf das Vorhandensein einer Unterschrift.

Hierzu wird ein prototypisches Evaluationsframework in Python entwickelt. Dieses führt die ausgewählten Verfahren auf einer gemeinsamen Datengrundlage aus, speichert die Extraktionsergebnisse und vergleicht sie mit manuell erstellten Referenzdaten. Die Untersuchung soll Unterschiede in der Extraktionsqualität sowie typische Fehlermuster sichtbar machen. Ergänzend werden Laufzeiten, Verarbeitungskosten und technische Anforderungen erfasst. Die erforderliche Auswahl und Aufbereitung relevanter \acs{pdf}-Seiten wird als vorgelagerter Verarbeitungsschritt berücksichtigt; der Schwerpunkt der Untersuchung liegt auf der Informationsextraktion und ihrer Evaluation.

Ein weiteres Ziel ist die wirtschaftliche Einordnung gegenüber der bisherigen manuellen Bearbeitung. Dazu sollen der Zeit- und Kostenaufwand beider Vorgehensweisen bis zu einem vergleichbaren Bearbeitungsstand gegenübergestellt werden. Beim automatisierten Vorgehen werden neben den Verarbeitungskosten auch der verbleibende menschliche Aufwand für Vorbereitung, Kontrolle und Korrektur berücksichtigt. Einmalige Entwicklungs- und Integrationsaufwände werden, soweit belastbar abschätzbar, gesondert ausgewiesen. Falls erforderliche Prozessdaten nicht vollständig erhoben werden können, erfolgt eine entsprechend gekennzeichnete Szenariobetrachtung auf Grundlage transparenter Annahmen.

Der Prototyp soll die extrahierten Angaben für eine menschliche Kontrolle und Korrektur bereitstellen. Eine produktive Anbindung an das betriebliche Schulungssystem gehört nicht zum verpflichtenden Implementierungsumfang. Die Möglichkeiten und Voraussetzungen einer solchen Anbindung werden bei der praktischen Einordnung betrachtet. Einsparungen durch einen bislang nicht umgesetzten automatisierten Import werden ausschließlich als bedingtes Szenario berücksichtigt. Die Arbeit soll auf dieser Grundlage eine nachvollziehbare Entscheidung über die Eignung der untersuchten Ansätze unterstützen. Ihre Aussagen beziehen sich auf die ausgewählten Modelle, die untersuchten Dokumente und die festgelegten Versuchsbedingungen.

\section{Forschungsfrage}

Aus der beschriebenen Problemstellung und Zielsetzung ergibt sich folgende zentrale Forschungsfrage:

\begin{quote}
  \researchquestion
\end{quote}

Zur systematischen Beantwortung wird die Forschungsfrage durch vier Unterfragen konkretisiert:

\begin{enumerate}
  \item Wie korrekt und vollständig werden die geforderten Informationen extrahiert und den jeweiligen Teilnehmenden zugeordnet, und welche Fehlermuster treten auf?
  \item Welcher Zeit- und Rechenaufwand entsteht bei der Verarbeitung durch die untersuchten Ansätze?
  \item Welche Kosten entstehen einschließlich manueller Kontrolle und Korrektur, und unter welchen Bedingungen ergeben sich Einsparungen gegenüber der bisherigen manuellen Bearbeitung?
  \item Welche Anforderungen und Einschränkungen ergeben sich hinsichtlich Dokumentvarianten, Datenschutz, technischer Integration und manueller Nachbearbeitung?
\end{enumerate}

Die Unterfragen verbinden die technische Leistungsbewertung mit den Bedingungen des betrieblichen Einsatzes. Die Wirtschaftlichkeitsbetrachtung konkretisiert dabei die Dimensionen Effizienz und praktische Einsatzfähigkeit. Angestrebt wird eine differenzierte Bewertung der Stärken, Schwächen und Zielkonflikte der untersuchten Ansätze. Ob ein Ansatz für den Anwendungsfall geeignet ist, soll anhand der erhobenen Ergebnisse und der betrieblichen Anforderungen beurteilt werden.

\section{Aufbau der Arbeit}

Kapitel~\ref{chap:grundlagen} erläutert die theoretischen Grundlagen der automatisierten Dokumentenverarbeitung und Informationsextraktion. Es behandelt \acs{ocr} und \ac{htr}, \acs{ki}-basierte Erkennungsansätze, Vision-Language-Modelle (\acs{vlm}) sowie hybride Verfahren. Darauf aufbauend werden Grundlagen der Bewertung und der relevante Forschungsstand dargestellt, um den Untersuchungsbedarf der Arbeit einzuordnen.

Kapitel~\ref{chap:methodik} beschreibt das methodische Vorgehen. Dazu gehören das Untersuchungsdesign, der Anwendungsfall und die Datengrundlage, die Erstellung der Ground Truth sowie die begründete Auswahl der Ansätze und Modelle. Außerdem werden Datenvorverarbeitung, Ausgabeformate, Evaluationsmetriken und die wirtschaftliche Bewertungsgrundlage festgelegt. Der Versuchsablauf und die Maßnahmen zur Reproduzierbarkeit schließen das Kapitel ab. Kapitel~\ref{chap:evaluationsframework} erläutert anschließend die Architektur des Evaluationsframeworks, die Implementierung der ausgewählten Ansätze sowie die automatisierte Durchführung und Auswertung der Versuche.

Kapitel~\ref{chap:ergebnisse} stellt die Ergebnisse zur Extraktionsqualität, zu Fehlermustern, Effizienz, Kosten und praktischer Einsatzfähigkeit dar und führt sie zu einer vergleichenden Gesamtbewertung zusammen. Kapitel~\ref{chap:diskussion} interpretiert die Befunde, ordnet sie in den Forschungsstand ein und beantwortet die Forschungsfrage. Dabei werden Zielkonflikte, betriebliche Einsatzbedingungen und Limitationen diskutiert. Kapitel~\ref{chap:fazit} fasst die zentralen Erkenntnisse zusammen, formuliert Schlussfolgerungen für den praktischen Einsatz und benennt weiterführende Untersuchungsmöglichkeiten.

% !TeX root = ../thesis.tex
% Abschnitte 2.1, 2.2, 2.4 und 2.5 ausgearbeitet; übrige Abschnitte bleiben im Entwurf.
\chapter{Theoretische Grundlagen und Stand der Forschung}
\label{chap:grundlagen}
Dieses Kapitel erläutert die Verarbeitung von Dokumentbildern, die Erkennung gedruckter und handschriftlicher Texte sowie Verfahren zur Gewinnung strukturierter Informationen. Die Begriffe und Verarbeitungsschritte bilden die Grundlage für den späteren Vergleich der Ansätze zur Extraktion aus Teilnehmerlisten.

\section{Automatisierte Dokumentenverarbeitung und Informationsextraktion}
Unter \ac{dia} werden Verfahren zur Erkennung und Auswertung visueller Bestandteile digitalisierter Dokumente zusammengefasst~\cite{OGorman1997DIA}. Bei den hier betrachteten \acs{pdf}-Dateien beginnt die Verarbeitung mit der Auswahl und Ausrichtung relevanter Seiten. Anschließend werden Bereiche lokalisiert, Texte erkannt und die Angaben den vorgesehenen Feldern zugeordnet. Diese Schritte können getrennt oder gemeinsam durch ein Modell umgesetzt werden.

Ein Teilbereich ist die \ac{dla}. O'Gorman und Kasturi unterscheiden dabei eine geometrische Ebene, auf der beispielsweise Textzeilen und Blöcke abgegrenzt werden, und eine funktionale Ebene, auf der die Bereiche eine Bedeutung erhalten~\cite{OGorman1997DIA}. Bei Teilnehmerlisten betrifft dies die Erkennung von Zeilen und Spalten sowie die Zuordnung zu Feldern wie Vorname, Personalnummer oder Unterschrift. Eine richtig transkribierte Zeichenfolge im falschen Feld ist deshalb kein korrektes Extraktionsergebnis. Für die Fehleranalyse ist diese Unterscheidung hilfreich, auch wenn ein Verfahren beide Aufgaben in einem Schritt bearbeitet.

Die Informationsextraktion überführt erkannte Inhalte in ein vorgegebenes Zielschema. Das Ergebnis soll für jede erfasste Teilnehmerzeile Vorname, Nachname, Personalnummer, Abteilung und das Vorhandensein einer Unterschrift enthalten. Dafür müssen zusammengehörige Angaben derselben Zeile zugeordnet werden. Skalický et al. unterscheiden die Extraktion einzelner Schlüsselinformationen von der Erkennung zusammengehöriger Einträge innerhalb eines Dokuments~\cite{Skalicky2022Benchmarks}. Fehlende und nicht lesbare Werte sind dabei gesondert zu kennzeichnen. In der späteren Bewertung werden Transkriptionsfehler, falsche Zeilen- oder Feldzuordnungen und unbegründete Ergänzungen getrennt betrachtet.

Bekannte Formularpositionen können die Lokalisierung bei vorlagenbasierten Verfahren vereinfachen. Bei festen Koordinaten müssen jedoch Verschiebungen, Skalierung und Drehung berücksichtigt werden. Layoutorientierte Modelle wie LayoutLM zeigen einen anderen Ansatz, indem sie Text und räumliche Position gemeinsam für Aufgaben des Dokumentenverständnisses modellieren~\cite{Xu2020LayoutLM}. Ob solche Ansätze für die vorliegenden handschriftlichen Listen geeignet sind, ist Gegenstand der Untersuchung.
\section{\texorpdfstring{\acs{ocr}}{OCR} und \texorpdfstring{\acs{htr}}{HTR}}
\acl{ocr} (\acs{ocr}) bezeichnet die Überführung von Textbildern in maschinenlesbaren Text. Der Begriff umfasst auch die Erkennung handschriftlicher Eingaben; \acs{trocr} wird beispielsweise als \acs{ocr}-Verfahren vorgestellt und auf gedruckten sowie handschriftlichen Textbildern untersucht~\cite{Li2023TrOCR}. \acl{htr} (\acs{htr}) bezeichnet in dieser Arbeit die speziell auf Handschrift bezogene Erkennungsaufgabe. Diese Arbeitsdefinition unterstützt die spätere Auswertung, ohne beide Begriffe als grundsätzlich getrennte Verfahrensklassen zu behandeln.

Handschriftliche Einträge können zwischen Schreibenden und innerhalb eines Dokuments in Form, Größe und Verbindung der Zeichen variieren. Besonders bei verbundener Schrift ist die Segmentierung in einzelne Zeichen schwierig; Erkennungsverfahren können daher auf Wort- oder Zeilenbildern ansetzen~\cite{Ploetz2009MarkovSurvey}. \ac{ctc} ermöglicht es Sequenzmodellen, Ausgabefolgen ohne vorgegebene Ausrichtung jedes Zeichens zum Eingabesignal zu lernen~\cite{Graves2006CTC}. Die Originalarbeit belegt damit keine konkrete Erkennungsleistung auf den hier untersuchten Teilnehmerlisten.

Je nach Dokument und Modell können vor der Erkennung Bildbereiche ausgerichtet, entrauscht, binarisiert oder segmentiert werden. Bei einer globalen Binarisierung gilt ein Schwellenwert für das gesamte Bild; Otsus Verfahren ermittelt ihn aus dem Grauwert-Histogramm~\cite{Otsu1979Threshold}. Lokale Verfahren wie das von Sauvola und Pietikäinen beschriebene berücksichtigen die Umgebung einzelner Bildbereiche und können dadurch auf wechselnde Bildbedingungen reagieren~\cite{Sauvola2000Binarization}. Ob eine Binarisierung hilfreich ist, muss für die eingesetzten Modelle geprüft werden.

Zur Korrektur einer Seitenschräglage beschreiben O'Gorman und Kasturi unter anderem Projektionsprofile und die Hough-Transformation~\cite{OGorman1997DIA}. Anschließend können Textzeilen oder Formularbereiche segmentiert werden. Bei handschriftlich ausgefüllten Teilnehmerlisten ist zu prüfen, wie sich berührte oder überschrittene Zellgrenzen auf diese Schritte auswirken. Auch nach einer korrekten Transkription müssen erkannte Angaben noch der richtigen Zeile und dem richtigen Feld zugeordnet werden.

Mehrstufige Systeme können Fehler früher Verarbeitungsschritte an spätere Schritte weitergeben. Daneben gibt es Ansätze, die Aufgaben zusammenfassen: \acs{trocr} konzentriert sich auf die Erkennung eines Textbilds~\cite{Li2023TrOCR}, während Donut ein Dokumentbild ohne vorgeschaltetes externes \acs{ocr}-System direkt in eine aufgabenspezifische Ausgabe überführt~\cite{Kim2022Donut}. Die Verfahren lösen damit unterschiedliche Teilaufgaben; ihre Leistung für die vorliegenden Listen muss unter vergleichbaren Bedingungen ermittelt werden.
\section{\texorpdfstring{\acs{ki}}{KI}- und Deep-Learning-Verfahren zur Handschrifterkennung}
Spezialisierte Erkennungsmodelle erhalten typischerweise ein Bild eines Textbereichs und geben eine Zeichen- oder Tokenfolge zurück. Eine solche Architektur kann die Transkription verbessern, übernimmt aber die Feldzuordnung nur dann, wenn diese ausdrücklich Teil der Eingabe oder eines zusätzlichen Verarbeitungsschritts ist. Für den Vergleich in dieser Arbeit ist deshalb die gesamte benötigte Pipeline zu betrachten und nicht allein das Erkennungsmodell.

\acs{trocr} ist ein Beispiel für einen Transformer-basierten Ansatz zur Texterkennung: Ein Bild-Transformer kodiert das Textbild, während ein Text-Transformer die Ausgabe erzeugt. Die Originalarbeit untersucht unter anderem handschriftliche Textbilder und beschreibt ein Vorgehen mit Vortraining und anschließendem Fine-Tuning~\cite{Li2023TrOCR}. Daraus folgt noch keine Eignung des Modells für die vorliegenden Teilnehmerlisten; insbesondere die Eingabeform, Sprache, Feldzuordnung und verfügbare Rechenumgebung müssen vor einer Auswahl geprüft werden.
\section{Vision-Language-Modelle}
Vision-Language-Modelle (\acs{vlm}) verarbeiten Bild- und Textinformationen gemeinsam. Für die Dokumentenextraktion ist besonders relevant, ob ein Modell gesuchte Angaben aus einem Dokumentbild oder Bildausschnitt in eine vorgegebene Struktur übertragen kann. Die Bezeichnung umfasst unterschiedliche Architekturen: Dokumentenspezifische Modelle wie Donut erzeugen aufgabenspezifische Ausgaben aus Dokumentbildern, während allgemeine, instruktionsgesteuerte Modelle Bildinformationen mit einem Sprachmodell verknüpfen~\cite{Kim2022Donut,Liu2023VisualInstruction}. Diese Gruppen sind nicht mit einem einheitlichen Leistungs- oder Ressourcenprofil gleichzusetzen.

\subsection{OCR-freies Dokumentenverständnis am Beispiel Donut}
Donut verzichtet bei der Anwendung auf ein vorgeschaltetes, separates \acs{ocr}-System. Ein visueller Transformer kodiert das Eingabebild; ein Textdecoder erzeugt daraus eine Folge von Ausgabetokens. Die Originalarbeit verwendet einen Swin-Transformer als visuellen Encoder und einen vortrainierten multilingualen Textdecoder. Aufgabenspezifische Prompt-Tokens steuern die Ausgabe, die anschließend in ein strukturiertes Format wie \ac{json} überführt werden kann~\cite{Kim2022Donut}. OCR-frei bedeutet in diesem Zusammenhang, dass kein eigenständiges OCR-Modul vorgeschaltet wird; das Modell muss den für die Aufgabe relevanten Text weiterhin aus dem Bild erschließen.

Die autoregressive Ausgabe lässt sich bei einem vorgegebenen Aufgabenprompt $q$ als bedingte Folge von Tokens beschreiben:
\begin{equation}
  P(Y\mid I,q)=\prod_{t=1}^{T}P(y_t\mid y_{<t},I,q).
  \label{eq:donut-tokenfolge}
\end{equation}
Dabei ist $I$ das Dokumentbild, $Y=(y_1,\ldots,y_T)$ die erzeugte Tokenfolge und $y_{<t}$ die vor dem Schritt $t$ erzeugte Teilfolge. Die Gleichung beschreibt die sequenzielle Generierung und legt weder die Feldqualität noch die Zuverlässigkeit der Ausgabe auf handschriftlichen Teilnehmerlisten fest~\cite{Kim2022Donut}.

\subsection{Allgemeine multimodale Modelle und Grenzen}
Allgemeine multimodale Modelle können durch visuelle Instruktionsabstimmung darauf trainiert werden, Bild- und Textanweisungen gemeinsam zu verarbeiten. \ac{llava} verbindet beispielsweise einen visuellen Encoder mit einem \ac{llm} und wird mit Bild-Text-Instruktionen angepasst~\cite{Liu2023VisualInstruction}. Dieser Nachweis allgemeiner visueller Instruktionsverarbeitung belegt für sich genommen noch keine zuverlässige Extraktion aus den vorliegenden Formularen.

OCRBench v2 zeigt bei den dort untersuchten multimodalen Modellen Grenzen unter anderem bei der Textlokalisierung, der Extraktion handschriftlicher Inhalte und dem Layoutverständnis~\cite{Fu2025OCRBenchV2}. Für Teilnehmerlisten sind deshalb neben der Zeichenfolge insbesondere die Zuordnung zu Zeilen und Feldern sowie fehlende, unlesbare oder unbegründet ergänzte Werte zu prüfen. Das Risiko plausibel wirkender, aber bildseitig nicht belegter Angaben wird in der Fehleranalyse untersucht; eine konkrete Häufigkeit lässt sich aus den veröffentlichten Benchmarks für diesen Anwendungsfall nicht ableiten. Ebenso sind Laufzeit, Speicherbedarf und Kosten für die tatsächlich ausgewählten Modelle unter den später festgelegten Bedingungen zu erfassen. Ergebnisse einer Kostenstudie zu gedruckten Berichten aus dem \acs{erp}-Bereich lassen sich nicht unmittelbar auf handschriftliche Listen übertragen~\cite{Lertyosbordin2026ERP}.
\section{Hybride Verfahren}
Hybride Verfahren verbinden mehrere spezialisierte Verarbeitungsschritte. Eine mögliche Pipeline für die vorliegende Aufgabe lokalisiert zunächst relevante Zeilen oder Felder, transkribiert die ausgeschnittenen Textbereiche und ordnet die Ergebnisse anschließend den gesuchten Attributen zu. Im Unterschied zu einem direkt generierenden Dokumentmodell bleiben Bildlokalisierung, Texterkennung und Strukturierung als einzelne Schritte erkennbar~\cite{Kim2022Donut}.

\subsection{Aufbau einer mehrstufigen Extraktionspipeline}
Die Lokalisierung liefert Bildausschnitte und gegebenenfalls Koordinaten der erkannten Bereiche. Ein darauf angewandtes \acs{ocr}- oder \acs{htr}-Modell erzeugt die Transkription. \acs{trocr} ist ein Beispiel für die Erkennung eines vorgegebenen Textbilds; die Ermittlung der Feldkoordinaten gehört nicht zu dieser Transkriptionsaufgabe und müsste separat erfolgen~\cite{Li2023TrOCR}. Für eine spätere Kontrolle sind Dokumentkennung, Bildbereich und erkannter Text gemeinsam zu speichern.

Im anschließenden Schritt können Regeln oder ein \acs{llm} die erkannten Textfragmente, ihre Positionen und gegebenenfalls Feldbezeichner in das Zielschema überführen. Bei Teilnehmerlisten betrifft dies besonders die Zuordnung mehrerer Werte zu derselben Person und die Unterscheidung zwischen fehlenden und unlesbaren Einträgen. Eine Sprachmodellkomponente ist dabei eine mögliche Ausgestaltung, keine notwendige Eigenschaft jeder hybriden Pipeline. Ihre Ausgabe muss anhand der Originalbilder und der Referenzdaten bewertet werden.

\subsection{Fehlerfortpflanzung, Nachprüfbarkeit und Aufwand}
Ein Fehler in der Lokalisierung oder Transkription kann nachfolgende Schritte beeinflussen. Kim et al. beschreiben diese Fehlerfortpflanzung für dokumentbezogene Extraktionspipelines~\cite{Kim2022Donut}. Van Strien et al. zeigen darüber hinaus für andere nachgelagerte Textverarbeitungsaufgaben, dass \acs{ocr}-Fehler deren Ergebnisse beeinträchtigen können~\cite{vanStrien2020OCRQuality}. Diese Befunde begründen eine getrennte Fehleranalyse, liefern aber keine Fehlerquote für die hier untersuchten Listen. Ohne Zugriff auf den betreffenden Bildbereich kann eine nachgelagerte Sprachmodellkomponente eine falsch erkannte Personalnummer nicht verlässlich aus dem Text allein rekonstruieren.

Gespeicherte Koordinaten ermöglichen es, extrahierte Angaben bei der menschlichen Prüfung zum Ursprungsbild zurückzuführen. Dieser Vorteil entsteht nur, wenn Lokalisierung und Feldzuordnung korrekt sind und die Verknüpfung in der Ergebnisstruktur erhalten bleibt; er ist nicht jeder hybriden Architektur automatisch gegeben~\cite{Skalicky2022Benchmarks}. Die getrennten Komponenten lassen sich einzeln prüfen und gegebenenfalls austauschen, verursachen aber auch eigenen Integrations- und Laufzeitaufwand. Ob die Pipeline gegenüber einem direkt generierenden Modell genauer, schneller oder kostengünstiger ist, wird anhand derselben Dokumente und eines vergleichbaren Bearbeitungsstands untersucht.
\section{Bewertung von Informationsextraktionssystemen}
\begin{inhalte}
  \item Extraktionsqualität
  \item Fehlermaße
  \item Laufzeit und Rechenaufwand
  \item Kosten
  \item Praktische Einsatzfähigkeit
\end{inhalte}
\section{Stand der Forschung und Forschungslücke}
\begin{inhalte}
  \item Bisherige Vergleiche von \acs{ocr}-, \acs{htr}- und multimodalen Modellen
  \item Offene Fragestellungen bei handschriftlichen strukturierten Dokumenten
  \item Einordnung der eigenen Untersuchung
\end{inhalte}

% !TeX root = ../thesis.tex
% Vorläufige Gliederung mit geplanten Inhalten.
\chapter{Methodik}
\label{chap:methodik}
\section{Untersuchungsdesign}
\begin{inhalte}
  \item Vergleich der untersuchten Ansätze unter kontrollierten Testbedingungen
  \item Reproduzierbare Versuchsdurchführung
\end{inhalte}
\section{Anwendungsfall und Datengrundlage}
Die Datengrundlage bilden historische, handschriftlich ausgefüllte Teilnehmerlisten aus Trainings und Schulungen eines Unternehmens. Der endgültige Umfang und die Zusammensetzung des Bestands stehen vor der Sichtung und Vorverarbeitung noch nicht fest. Für die Stichprobenbeschreibung werden später insbesondere die Zahl der Dokumente und Teilnehmerzeilen, die Dokumentvarianten, die Bildqualität und die Häufigkeit fehlender oder nicht lesbarer Einträge erfasst. Erst auf dieser Grundlage können Auswahl und Aussagebereich der Untersuchung abschließend begründet werden.

Als Zielfelder sind Vorname, Nachname, Personalnummer und Abteilung vorgesehen. Beim Unterschriftsfeld wird ausschließlich erfasst, ob eine Unterschrift vorhanden ist. Eine inhaltliche Analyse der Unterschrift oder Identitätsprüfung ist nicht Gegenstand der Untersuchung. Ob die Personalnummer als eindeutige Kennung für die Aufteilung wiederkehrender Personen genutzt werden kann, wird bei der Datensichtung geprüft.

Die Listen enthalten personenbezogene Angaben. Eine Übermittlung von Dokumenten oder Ausschnitten an externe \acs{api}-basierte Modelldienste kommt nur nach Freigabe oder nach geprüfter Anonymisierung in Betracht. Bis zur Klärung dieser Voraussetzung bleibt die Verarbeitung der Originaldaten lokal. Dieser Rahmen beeinflusst die spätere Modellauswahl und ist im Versuchsprotokoll zu dokumentieren.
\section{Erstellung der Ground Truth}
Die Ground Truth wird unabhängig von den Ausgaben der untersuchten Modelle erstellt oder manuell überprüft. Als Annotationseinheit dient zunächst die Teilnehmerzeile mit den zugehörigen Zielfeldern. Die Dokumentkennung und eine stabile Zeilenkennung ermöglichen die spätere Zuordnung zwischen Bild, Referenz und Modellausgabe. Vor Beginn der Annotation wird ein Codebuch mit Beispielen für alle Feldtypen festgelegt.

Das Codebuch unterscheidet einen tatsächlich fehlenden Eintrag von einem vorhandenen, aber nicht zuverlässig lesbaren Eintrag. Für unklare Fälle wird die Entscheidung dokumentiert und möglichst durch eine zweite Sichtung überprüft. Die originale Ablesung bleibt erhalten; eine gegebenenfalls vereinheitlichte Schreibweise wird separat gespeichert. Bei Namen, Abteilungen und Personalnummern müssen insbesondere Groß- und Kleinschreibung, Leerzeichen, Sonderzeichen und führende Nullen vor der Auswertung geregelt werden. Diese Regeln dürfen nach Sichtung der finalen Modellergebnisse nicht mehr geändert werden.

Für die Auswertung werden nur Felder verwendet, deren Referenzwert nach dem festgelegten Protokoll bestimmbar ist. Die Zahl ausgeschlossener oder als unleserlich markierter Felder wird berichtet. Für das Unterschriftsfeld wird die Anwesenheit einer Unterschrift binär annotiert. Die Qualitätskontrolle der Referenzdaten wird nach Sichtung des Bestands festgelegt und begründet.
\section{Auswahl der untersuchten Ansätze und Modelle}
Für den Vergleich werden wenige technisch unterschiedliche, im Projektumfang ausführbare Ansätze ausgewählt. Als Klassen kommen spezialisierte \acs{ocr}-/\acs{htr}-Modelle, dokumentenspezifische Modelle, allgemeine Vision-Language-Modelle und gegebenenfalls eine hybride Pipeline infrage. Die in Kapitel~\ref{chap:grundlagen} genannten Modelle sind Literaturbeispiele und noch keine endgültige Versuchsauswahl.

Die Auswahl wird vor der finalen Evaluation protokolliert. Zu prüfen sind die Eignung für handschriftliche Eingaben und die benötigte Eingabeeinheit, die Möglichkeit einer vollständigen Feldextraktion, die Verfügbarkeit einer dokumentierten Modellversion, der Aufwand der lokalen Ausführung beziehungsweise die Zulässigkeit einer Nutzung einer \acs{api} sowie die verfügbare Hardware. Eine zusätzliche hybride Pipeline wird nur aufgenommen, wenn sie gegenüber den Einzelverfahren eine klar prüfbare Fragestellung liefert und der Vergleich mit denselben Dokumenten möglich bleibt.
\section{Datenvorverarbeitung}
\begin{inhalte}
  \item Auswahl relevanter Seiten aus den \acs{pdf}-Dateien
  \item Vereinheitlichung der Seitenorientierung
  \item Bildvorverarbeitung
  \item Gegebenenfalls Segmentierung beziehungsweise Feldextraktion
\end{inhalte}
\section{Einheitliche Extraktions- und Ausgabeformate}
\begin{inhalte}
  \item Einheitliche Eingabe- und Ausgabeformate
  \item Gewünschte strukturierte Ausgabe
\end{inhalte}
\section{Evaluationsmetriken}
\subsection{Extraktionsqualität}
\begin{inhalte}
  \item Exact Match
  \item \ac{cer}
  \item Word Error Rate
  \item Gegebenenfalls Precision, Recall und F1
\end{inhalte}
\subsection{Effizienz}
\begin{inhalte}
  \item Laufzeit
  \item Rechenaufwand
\end{inhalte}
\subsection{Kosten}
\begin{inhalte}
  \item Inferenz- beziehungsweise \acs{api}-Kosten
  \item Zeit- und Kostenaufwand der manuellen Erfassung sowie der Vorbereitung, Kontrolle und Korrektur automatisierter Ergebnisse bis zu einem vergleichbaren Bearbeitungsstand
  \item Einmaliger Entwicklungs- und Integrationsaufwand getrennt von laufenden Kosten
  \item Bei fehlenden Prozessdaten: gekennzeichnete Szenarien mit transparenten Annahmen
\end{inhalte}
\subsection{Praktische Einsatzfähigkeit}
\begin{inhalte}
  \item Datenschutz, technische Integration und manueller Nachbearbeitungsaufwand
  \item Voraussetzungen einer möglichen Anbindung an das Schulungssystem
\end{inhalte}
\section{Versuchsablauf und Reproduzierbarkeit}
Nach der Bestandsaufnahme werden Entwicklungs- und Testdokumente vor der Anpassung von Prompts, Parametern oder Nachverarbeitungsregeln getrennt. Dokumente aus dem finalen Testteil werden für diese Anpassungen nicht verwendet. Die konkrete Aufteilung wird anhand des tatsächlichen Bestands festgelegt und vor den Testläufen dokumentiert. Falls dieselben Personen oder sehr ähnliche Listen mehrfach vorkommen, muss die Aufteilung mögliche Überschneidungen berücksichtigen.

Für jeden Lauf werden Dokumentkennung, Modell- und Versionsangaben, Eingabevorbereitung, Prompt beziehungsweise Parameter, unbearbeitete Modellausgabe, normalisierte Ausgabe, Laufzeit und gegebenenfalls Kosten gespeichert. Erst danach erfolgt der Abgleich mit der Ground Truth. Fehler werden getrennt nach Transkription, Feldzuordnung, Auslassung und unbegründeter Ergänzung analysiert. Für nicht deterministische Verfahren wird vor der finalen Evaluation festgelegt, ob Wiederholungsläufe erforderlich sind; ihre Ergebnisse dürfen nicht stillschweigend zu einem einzigen Lauf zusammengefasst werden.

% !TeX root = ../thesis.tex
% Vorläufige Gliederung mit geplanten Inhalten.
\chapter{Technische Umsetzung des Evaluationsframeworks}
\label{chap:evaluationsframework}
\section{Aufbau und Systemarchitektur}
\begin{inhalte}
  \item Modularer Aufbau
  \item Einheitliche Modell- beziehungsweise Pipeline-Schnittstellen
  \item Reproduzierbare Verarbeitung
  \item Einlesen der Dokumente und Bildaufbereitung
  \item Gegebenenfalls Feld- oder Bereichserkennung
\end{inhalte}
\section{Implementierung der untersuchten Ansätze}
\begin{inhalte}
  \item Integration spezialisierter \acs{ocr}-/\acs{htr}- beziehungsweise Machine-Learning-Ansätze
  \item Integration von Vision-Language-Modellen: Prompting, Bildverarbeitung, strukturierte Extraktion und Ausgabe im Format \ac{json}
  \item Kombination mehrerer Modelle beziehungsweise Verarbeitungsschritte
  \item Beispielsweise \acs{ocr}/\acs{htr} mit anschließender Strukturierung oder Korrektur durch ein \ac{llm}
  \item Standardisierte Ergebnisstruktur
\end{inhalte}
\section{Automatisierte Durchführung und Evaluation}
\begin{inhalte}
  \item Automatisierter Ground-Truth-Vergleich
  \item Speicherung von Modellantworten und Metriken
  \item Erfassung von Laufzeit und Kosten
  \item Gegebenenfalls Kennzeichnung unsicherer Extraktionen anhand einfacher Konsistenz- oder Confidence-Verfahren
\end{inhalte}

% !TeX root = ../thesis.tex
% Vorläufige Gliederung mit geplanten Inhalten.
\chapter{Ergebnisse und Analyse}
\label{chap:ergebnisse}
\section{Ergebnisse der Extraktionsqualität}
\begin{inhalte}
  \item Vergleich der untersuchten Ansätze
  \item Auswertung der zentralen Qualitätsmetriken
\end{inhalte}
\section{Analyse der untersuchten Ansätze}
\begin{inhalte}
  \item Gegenüberstellung der untersuchten Modellklassen und Verarbeitungspipelines
  \item Unterschiede zwischen spezialisierten Modellen, Vision-Language-Modellen und hybriden Ansätzen
\end{inhalte}
\section{Fehleranalyse}
\begin{inhalte}
  \item Typische Fehlerarten
  \item Unterschiede zwischen Modellklassen
  \item Einfluss von Handschrift und Dokumentqualität
\end{inhalte}
\section{Analyse der Effizienz}
\begin{inhalte}
  \item Laufzeit
  \item Rechenaufwand
  \item Anzahl beziehungsweise Komplexität der Inferenzschritte
\end{inhalte}
\section{Kostenanalyse}
\begin{inhalte}
  \item \acs{api}-Kosten
  \item Lokale Inferenz
  \item Kosten pro Dokument beziehungsweise Extraktion
  \item Manueller Aufwand für Vorbereitung, Kontrolle und Korrektur
  \item Vergleich mit der bisherigen manuellen Erfassung bei gleichem Bearbeitungsstand
  \item Einmalige Aufwände und bedingte Szenarien getrennt ausweisen
\end{inhalte}
\section{Analyse der praktischen Einsatzfähigkeit}
\begin{inhalte}
  \item Technische Umsetzbarkeit
  \item Integrationsaufwand
  \item Hardware- und Infrastrukturbedarf
  \item Datenschutz und Informationssicherheit
  \item Skalierbarkeit und Automatisierbarkeit im Unternehmenskontext
\end{inhalte}
\section{Gesamtbewertung der untersuchten Ansätze}
\begin{inhalte}
  \item Gegenüberstellung der Ergebnisse zu Qualität, Effizienz, Kosten und praktischer Einsatzfähigkeit
\end{inhalte}

% !TeX root = ../thesis.tex
% Vorläufige Gliederung mit geplanten Inhalten.
\chapter{Diskussion}
\label{chap:diskussion}
\section{Interpretation und Einordnung der Ergebnisse}
\begin{inhalte}
  \item Interpretation der Ergebnisse vor dem Hintergrund der Forschungsfrage
  \item Einordnung in den Stand der Forschung
\end{inhalte}
\section{Beantwortung der Forschungsfrage}
\begin{inhalte}
  \item Synthese und Beantwortung der Leitfrage: \textit{\researchquestion}
\end{inhalte}
\section{Trade-offs zwischen den untersuchten Ansätzen}
\begin{inhalte}
  \item Trade-offs zwischen Qualität, Effizienz, Kosten und Implementierungskomplexität
\end{inhalte}
\section{Praktische Einordnung im Unternehmenskontext}
\begin{inhalte}
  \item Eignung der Ansätze für betriebliche Prozesse
  \item Technische und organisatorische Anforderungen an den Einsatz
\end{inhalte}
\section{Limitationen der Untersuchung}
\begin{inhalte}
  \item Datengrundlage
  \item Modellauswahl
  \item Generalisierbarkeit
  \item Technische und organisatorische Einschränkungen
\end{inhalte}

% !TeX root = ../thesis.tex
% Vorläufige Gliederung mit geplanten Inhalten.
\chapter{Fazit und Ausblick}
\label{chap:fazit}
\section{Zusammenfassung der zentralen Ergebnisse}
\section{Schlussfolgerungen für den praktischen Einsatz}
\begin{inhalte}
  \item Abschließende Beurteilung der Eignung spezialisierter Modelle, Vision-Language-Modelle und hybrider Ansätze
\end{inhalte}
\section{Ausblick}
\begin{inhalte}
  \item Mögliche Weiterentwicklung des Forschungsprototyps
  \item Weitere Modelle und Datensätze
  \item Mögliche spätere Integration in Unternehmensprozesse
\end{inhalte}


% references.bib wird automatisch durch Zotero + Better BibTeX gepflegt.
% verified_sources.bib enthält bereits geprüfte Primärquellen des Entwurfs.
% Nur tatsächlich zitierte Quellen werden ausgegeben.
\bibliographystyle{IEEEtran}
\bibliography{references,verified_sources}

\appendix
% !TeX root = ../thesis.tex
\chapter{Ergänzende Materialien}
\label{app:materialien}

\TODO{Nur Material aufnehmen, das für die Nachvollziehbarkeit wichtig ist, aber den Lesefluss im Hauptteil stören würde. Wird kein Anhang benötigt, die zugehörige Zeile in \texttt{thesis.tex} entfernen.}


\end{document}
